\chapter[Sugar, Slavery and Empire]{How a Lump of Sugar Connects Slavery, Consumption and Empire}
\section{A Sugar Cube's Previous Life}
Pick up something from the table: a sugar cube.

It is small, about four grams, uniformly white and neatly edged. Drop it into milk tea or coffee and it vanishes within seconds. It is so cheap you have never given it a moment's thought: a whole convenience-store box costs less than one milk tea. Among our age's most unremarkable seasonings, it often occupies a supermarket's lowest shelf.

Now try a thought experiment: give this cube to someone three hundred years ago.

Give it to a fashionable London lady in 1700 and she will lift it with silver tongs and break off a little. Sugar was then a luxury priced in small pieces, locked in a special box whose key hung at the mistress's waist. Give it to an enslaved African on a contemporary Barbados plantation and he will recognise it: the cane he cut, the mill he pushed and the copper pan he watched all night eventually became these white crystals. Yet he may never have tasted refined sugar in his life. Give it to a Canton tea merchant in 1750 and he will say it is nothing new. Chinese people have eaten sugar for over a millennium, mostly from southern cane fields, by another route.

One sugar cube becomes three different things in three hands: a status symbol, crystallised suffering and an everyday commodity. How did a single chain bind these identities together?

A clue lies in the word itself. English ``sugar'', French ``sucre'' and Spanish ``azúcar'' derive from Arabic ``sukkar'', borrowed from Persian ``shakar'', which goes back to Sanskrit, ancient India's language: ``śarkarā''. Surprisingly, its earliest meaning was gravel or little stones, since early coarse sugar crystals resembled sand. A word's journey follows a commodity's. Sugar cane was first domesticated around New Guinea, then spread west to India, where people first mastered boiling its juice into crystals. Sugar continued west with Persians, Arab traders and Islamic expansion, reaching the Mediterranean.

China had its own stop on this route. The New Book of Tang's account of the Western Regions records that during the Zhenguan era, Emperor Taizong sent envoys to Magadha, in today's north-eastern India, to learn sugar-making. Returning, they applied the method to tribute cane from Yangzhou, producing sugar whose colour and flavour ``far surpassed that of the Western Regions''. By Song times, Suining in Sichuan made frost-white rock sugar. Wang Zhuo devoted a book to it, the Treatise on Sugar Frost, the world's first specialised work on sugar manufacture.

Notice: Europeans neither invented sugar nor introduced it to the world. The story is the reverse. At a particular moment, Europe took an ancient Asian crop across the ocean and organised its production in an unprecedented way. That organisation was the plantation, and its labour force consisted of slaves.

How can four grams bear such historical weight? Begin with a scene.

\section{A Night in Barbados}
Enter the 1660s on Barbados, a small island at the Caribbean's eastern edge. It is an ordinary night in the cane-cutting season.

The signal ending daytime work sounded long ago, but the plantation shows no sign of stopping. Lanterns illuminate the mill house. Oxen turn three upright rollers with a dull creak. Enslaved Africans feed in armfuls of cut cane; another worker pulls the crushed stalks out opposite. This is the island's most dangerous job. A caught sleeve or finger can drag a whole body into the rollers in seconds, so an axe is always kept beside the mill. That is not my embellishment: contemporary plantation manuals prescribed it as operating procedure.

Juice runs through wooden channels into the adjoining boiling house, where another ordeal awaits. Copper pans, decreasing in size, stand over fires. Juice passes between them; scum must be constantly skimmed, and the sugar master's eye alone judges the heat. The room is a steamer, its air thick with the sickening, almost rank sweetness of molasses. Exhausted night workers doze standing up. During harvest, owners wish for forty-eight-hour days: cut cane spoils unless promptly crushed, and juice sours unless promptly boiled. People can be replaced; the season will not wait.

A few kilometres away lies the island's other face. Only twenty or thirty years earlier, Barbados was forested. Now cane has almost devoured the trees. Timber, food and even vegetables must be imported. More than half the inhabitants are enslaved Africans, a proportion still rising. Small English farmers, unable to profit from tobacco, sell up and leave. Those remaining are large planters who can afford mills, copper pans and slaves. Historians call this the ``sugar revolution'': one crop overturning an island's society, ecology and population.

Cut to London, more than six thousand kilometres away, on the same night---morning there. A merchant household enriched by overseas trade breakfasts on white bread, butter and tea. The younger daughter stretches up to steal from the sugar bowl. Her mother taps her hand, not worrying about the expense, but her teeth. Yesterday, at a coffeehouse, the father concluded a deal underwriting a cargo ship bound for the West Indies.

An Atlantic separates the scenes. These people will never meet and may not know the others exist. Yet the sweetness on the London girl's fingertips and the axe beside the Barbados mill belong to opposite ends of one chain.

\section{What Lies at Sweetness's Other End?}
Set out the conflict before reaching conclusions.

Two accounts of sugar both sound forceful.

The first: sugar is merely a crop and commodity, innocent in itself. For tens of thousands of years, humans have domesticated plants and improved food. Cane makes bitterness sweet and blandness pleasant. What is wrong with that? The London girl should not answer for an axe she does not know exists. If consumption itself is culpable, none of us escapes: who knows what lies deep in the chains behind your shoes or phone? Is linking sweetness to violence a form of moral blackmail?

The second: sugar was never merely a commodity. In the seventeenth and eighteenth centuries, it was among the world's most profitable crops, its profits squeezed not from soil but from human bodies. Mills crushed cane; plantations crushed people. The Atlantic system built for sugar, and later cotton and coffee, packed over twelve million Africans into ships bound for America. More than a million died at sea. This was no residue of an ancient institution: it was one of the world's most modern industries, with international finance, marine insurance, transnational supply chains and precise cost accounting. Calling it merely business resembles calling fire merely light. It illuminated some drawing rooms by consuming other people's entire lives.

Each account holds part of the facts. Their collision brings us to the real question:

\textbf{When an object sits before you, how far are you from the violence behind it? Can that distance release you from responsibility?}

This is difficult because it asks neither historical knowledge---archives can establish what happened in Barbados---nor a moral slogan: slavery was plainly wrong. It poses a structural problem. In a world of elaborate specialisation and long chains, enjoyment and cost are deliberately placed at opposite ends. Those tasting sweetness cannot see suffering. Who designed this? Are uninformed consumers innocent, or is their ignorance itself part of the chain?

We need more components before answering. First hear what people along the chain said for themselves.

\section[Four Voices along the Chain]{Four Voices: Accounts, Ship Holds, Tea Tables and Parliament}
Slavery was not a vague cloud of past evil, but concrete daily decisions. We need at least four perspectives to understand the machine.

\textbf{Voice one: the planter's ledger.}

Enter an eighteenth-century Jamaican planter's counting house. The first object is a ledger, not a whip. Plantations demanded enormous capital: land, mills, pans and slaves required an initial investment comparable to an English factory. The accountant calculated precisely. European indentured labourers needed food and lodging for several years, then departed free, with land to be allocated. Disease, enslavement and massacre had nearly exterminated the island's Indigenous people within decades of Columbus, leaving nobody to hire. An African purchased at prevailing ``market rates'', however, could ``repay the investment'' within a few years. The ledger treated slaves not as people but as self-reproducing equipment subject to rapid wear. Caribbean plantation deaths long exceeded births, requiring constant purchases to maintain numbers. The ledger had a special column for ``depreciation''.

Hear their defence fully, not to forgive them, but to understand the machine. Broadly: British, Spanish and French law permitted the trade; if we abstain, Dutch or French traders will profit instead; we rescue Africans from ``savagery'' and baptise them; West Indian slaves may eat no worse than Manchester workers. Uncomfortable as these sound now, each could win parliamentary gentlemen's approval. The last is revealing. Comparing plantation misery with factory misery exposes the defence's weakness: it does not establish slavery's rightness, only others' wrongness. A system reduced to ``everyone is dirty'' as its defence already knows something about itself.

\textbf{Voice two: enslaved Africans.}

First correct an easily mistaken picture: white men carrying guns deep into Africa and seizing everyone they saw. That is largely wrong. Europeans built dozens of coastal forts and trading stations, but for centuries rarely entered the interior, constrained by tropical disease. Most captives entered the trade through local warfare, debt, kidnapping or intercommunity judicial penalties, then passed through African merchants to European captains at coastal forts. This does not diminish European responsibility. It reveals its division: European purchasing power and American labour demand created a permanent market for African raiding, turning existing scattered enslavement into a production-line industry. People made choices at every link; profits ultimately pointed towards plantations across the sea.

The other side matters equally: captives were never silent cargo. They jumped overboard, refused food and mutinied. On plantations, they slowed work, feigned illness, damaged tools, poisoned owners' animals and, most radically, escaped. In Brazil, fugitives formed the inland federation of Palmares, surviving nearly a century and perhaps numbering tens of thousands at its height. Portuguese colonists required more than a dozen campaigns to destroy it. Jamaican Maroons fought so effectively from mountain refuges that Britain signed a treaty recognising their autonomy in 1739. Slaveholders negotiating with fugitives was inconceivable in a worldview of human cargo. The most dramatic chapter opened in 1791. French Saint-Domingue, now Haiti, was the world's largest sugar producer, supplying roughly forty per cent of global output. Hundreds of thousands rose, fought for over a decade, defeated Spanish and British forces and Napoleon's French expedition, and founded the Republic of Haiti in 1804. This was history's only slave uprising directly establishing a state. Anyone claiming slaves simply awaited liberation should read this page first.

Amid the violence, an entirely new society was also growing. Plantations compressed Africans, Europeans and surviving Indigenous people onto the same land. Over centuries they mixed, intermarried and learned from one another, creating unprecedented things: creole languages joining African vocabulary to European grammar; religions such as Candomblé combining African beliefs and Catholicism; predecessors of rumba and reggae; and cooking that placed African okra, Indigenous cassava and European cured meat in one pot. None was a copy of an older civilisation. These were New World creations born of Atlantic violence and creativity together. The character of cities from Havana to Rio de Janeiro is rooted in this history.

\textbf{Voice three: consumers at British tea tables.}

The chain's final link was not uniform either. By the late eighteenth century, sugar had become an everyday working-class food rather than a lady's luxury; we shall return to that transition. At its greatest popularity, some consumers questioned its origins. In 1787, potter Josiah Wedgwood made a medallion showing a kneeling, chained slave raising his hands, surrounded by ``Am I not a man and a brother?'' It became immensely popular on brooches, snuffboxes and ornaments, worn by countless people. After an abolition bill failed in Parliament in 1791, a spontaneous boycott swept Britain. Contemporary estimates put participating households in the hundreds of thousands, refusing West Indian sugar. William Fox's pamphlet declared, in effect, that every pound of West Indian sugar consumed meant eating two ounces of human flesh. Sugar bowls appeared marked ``East India Sugar---Not Made by Slaves''. Voting with your wallet was no internet-age invention. British housewives knew how two centuries ago.

\textbf{Voice four: parliamentary calculations.}

In Parliament, where decisions were made, arguments were cooler and colder. Anti-abolition MPs calculated that West Indian sugar was a major tax source; ending the slave trade would surrender markets to France; Liverpool and Bristol, enriched by slaves and sugar, would decline; British sea power rested on transoceanic shipping and sailors, so removing its largest route meant cutting off an arm. These calculations were real, explaining abolition's difficulty. Slavery was not only a moral evil, but a structure of interests entangled with imperial revenue, finance and local employment. Understanding this explains both the decades-long struggle and why abolition eventually happened.

\section[Thinking Bridge: Trace the Commodity]{Thinking Bridge: Following a Commodity Chain Backwards}
We now have a portable tool: \textbf{commodity-chain analysis}. It is as straightforward as unpacking a parcel. Pick any product, trace it backwards, draw its links and ask the same questions at each.

For your sweetened milk tea, the backward route looks roughly like this:

\begin{quote}
Consumption (you) $\rightarrow$ retail (tea shop, supermarket) $\rightarrow$ refining (sugar refinery) $\rightarrow$ long-distance transport (ships, insurance, finance) $\rightarrow$ initial processing (plantation boiling house) $\rightarrow$ cultivation and harvesting (cane fields) $\rightarrow$ origins of land and labour (whose land, whose hands?)
\end{quote}
At each link ask: \textbf{who works here, who profits here, and who bears the cost?}

Simple questions can be powerful because they treat a modern condition: ``chain myopia''. Longer chains make the beginning harder to see from the end. London's girl cannot see the Barbados mill, just as you cannot see the cobalt mine behind your phone battery. This is not entirely consumers' fault. Designers of the chain want its ends out of sight of one another: distance cushions conscience. Commodity-chain analysis reconnects what distance hides.

It also explains the famous textbook diagram of ``triangular trade''. Europe sent guns, cloth and rum to African coasts for slaves; captives crossed the Atlantic to American plantations; sugar, molasses, cotton and tobacco returned to Europe. Three sides make a triangle. Treat it as a map, not a photograph. Actual routes were far messier: New England fishing boats sold salted fish to feed Caribbean slaves; Brazil imported slaves directly for its gold mines; African merchants' inland routes formed separate networks. A simplified diagram reveals structure at a glance, but compresses people into arrows. Commodity-chain analysis reverses that move, placing people at both ends of every arrow.

The tool works on any product: a pencil, with timber, graphite, rubber and brass from different continents; a T-shirt; a phone. Even abstractions can be unpacked: an internet catchphrase---who created it, who gained traffic, who was mocked?---or a trending topic. Anything passing through many hands before reaching you deserves the question: where was its previous stop?

Practise at dinner tonight. Choose one ingredient, even a spoonful of sugar, and name five possible stops before your bowl. You may find your supposed common knowledge runs out at the second. That is precisely the chain's intended effect.

\section{A Survivor's Testimony}
Read a short primary source by someone who survived the chain's darkest link and learned to write in his oppressors' language.

According to his account, Olaudah Equiano was born around 1745 in present-day Nigeria, kidnapped at about eleven and passed between owners before being sold onto a transatlantic slave ship. He later moved between the Royal Navy, merchant vessels and plantations, purchased his freedom with forty saved pounds in 1766, settled in London and became a powerful abolitionist speaker. In 1789 he published The Interesting Narrative of the Life of Olaudah Equiano, or Gustavus Vassa, the African, Written by Himself. Nine editions appeared during his lifetime, with translations into several languages. Most of the twelve million transported captives left not a written word; this book therefore bears unusual weight.

In chapter two, he recalls the slave ship's hold. The following passages translate the Chinese renderings made by this book's author:

\begin{quote}
``Even while the ship remained near the coast, the hold's stench was unbearable, making a short stay dangerous \ldots{} Now that the whole `cargo' was confined together, the place became a seat of pestilence.''
\end{quote}
\begin{quote}
``The lack of air below, the extreme heat, and the number aboard---each person scarcely had space to turn---almost suffocated us.''
\end{quote}
He goes on to describe fever and dysentery spreading through the hold, and his wish to die. Pause at ``cargo'': irony and accusation together. In a world ruled by ledgers, people were goods. He deliberately adopts the cargo's voice to remind you that cargo speaks.

With evidence in hand, historians' next move is not emotion but questioning: is this testimony reliable, and how should we use it?

An honest difficulty must be disclosed. In 2005, scholar Vincent Carretta published research suggesting that surviving baptismal and naval records place Equiano's birth in South Carolina rather than Africa. If correct, the African childhood and Middle Passage descriptions may combine other survivors' accounts rather than personal experience. Scholars remain divided.

This controversy makes an excellent lesson in reading evidence. Separate three levels. First, records support that Equiano was enslaved, bought his freedom and became an abolitionist: these facts are undisputed. Second, independent evidence---captains' journals, merchants' accounts and later parliamentary testimony---corroborates crowding, stench and mortality below deck. Even if he did not personally experience the Middle Passage, the description's historical core remains sound. Third, a book written for abolition naturally seeks persuasion. That does not make it a lie, but demands the approach appropriate to any action-oriented text: compare sources and distinguish eyewitness experience, reported material and rhetoric.

That is basic primary-source literacy. Do not abandon verification because a voice belongs to a victim, or discard a whole testimony because one point is doubtful. Evidence acquires weight through corroboration with other evidence.

\section{How Were the Chains Unlocked?}
In 1833, Parliament passed the Slavery Abolition Act, effective the following year. About 800,000 enslaved people in the empire gained freedom, with a condition astonishing to modern readers: the government paid twenty million pounds, roughly forty per cent of annual expenditure, not to the enslaved but to \textbf{slaveholders}, compensating their ``property losses''. Borrowed through government debt, this money was not fully repaid by British taxpayers until 2015. The freed people received nothing and had to work four more unpaid years as ``apprentices'' for former masters.

Those are the facts. Now enter the arena: why did abolition occur? Keep four categories distinct. What happened---legislation, compensation, dates---is evidence. Why it happened is explanation, where three accounts will compete. Contemporary understandings---abolitionists citing conscience, opponents national ruin---are participants' statements. Our judgement that compensation went to the wrong recipients is evaluation. Confusing them makes historical mush.

\textbf{Explanation one: conscience won.}

The abolition movement was indeed humanity's first large-scale modern social movement. Quakers first prohibited slave trading among their members; Wilberforce introduced parliamentary proposals year after year; hundreds of thousands of Fox's pamphlets circulated; households boycotted sugar; petitions flooded Parliament. The evidence is strong: without this movement, the legislative timetable makes no sense. Its weakness is clear too. Why did British consciences awaken in the late eighteenth century? Christianity had existed alongside slavery for seventeen hundred years. Scripture had not changed: what had? Moral explanations identify the driving force better than the timing.

\textbf{Explanation two: calculation won.}

In 1944, Trinidadian scholar Eric Williams, later his country's founding prime minister, published the explosive Capitalism and Slavery. His argument had two steps. First, slave-trade and plantation profits nourished British banking, insurance and industrial capital: blood stained industrialisation's seed money. Second, by the late eighteenth century, exhausted soils and competition weakened West Indian plantations, turning slavery from asset to burden while Indian markets and Asian trade promised more. Britain therefore ``rationally'' discarded it. This elegantly fills the moral account's gap, explaining timing and the enthusiasm of new industrial interests already moving beyond the Caribbean. Critics, however, seized its limits. Later economic historians calculated that slavery remained profitable before abolition: the West Indies had not declined beyond viability. Nor does this explain hundreds of thousands of boycotting households, which bore costs without financial gain.

\textbf{Explanation three: enslaved people broke their own chains.}

Turn back to the Caribbean. News of Haiti's 1791 revolution reached plantations and parliaments alike. Slaveholders slept uneasily, and repression costs rose. Jamaican revolts followed one another. At Christmas 1831, sixty thousand slaves led by Samuel Sharpe launched the Baptist War, burning more than a hundred plantations. The revolt was brutally suppressed and Sharpe hanged; less than two years later, abolition passed. Here freedom was no gift from London drawing rooms, but a price generations of rebels raised with their lives. Slavery's maintenance costs exceeded its returns because enslaved people raised them. This account's great achievement is restoring them as historical agents rather than objects awaiting rescue. Its limitation: revolts alone could not overturn an empire's legal system. Haiti fought its freedom into being; British colonies ultimately received abolition through parliamentary votes. Weapons and debating chambers each contributed. Moreover, Haiti frightened some wavering slaveholders into hardened resistance. Rebellion advanced abolition but temporarily reinforced chains too.

Each explanation holds part, not all, of the truth. This is not compromise for its own sake; their respective functions are clear. Resistance raised slavery's costs, economic change altered imperial calculations, and moral activism translated these into parliamentary language. History rarely moves through one hand. Seeing each is more useful than arguing which mattered most.

\section[Further Challenge: Sweetness and Power]{Further Challenge: Sweetness and Power---The Anthropology of Sugar}
Now introduce the chapter's weightiest book. In 1985, American anthropologist Sidney Mintz published Sweetness and Power: The Place of Sugar in Modern History. Studying one commodity, it became widely recognised among the twentieth century's most important anthropological works by answering a question nobody had seriously asked: why sugar?

Mintz's central finding emerges from a timeline. Sugar entered the British diet along a clear downward social curve. First it was medicine, sold by apothecaries in tiny measured quantities for coughs and melancholy. Then came spice and luxury: aristocratic feasts displayed sugar castles, swans and entire fleets, declaring not ``eat this'' but ``I can afford to squander this''. Next it sweetened three new colonial drinks: tea, coffee and cocoa. Finally, in the nineteenth century, it became a staple. Factory workers breakfasted on white bread and jam, and took strong sweet tea in the afternoon. Over a century, annual British consumption rose from under two kilograms per person to around eight, fastest among the poorest. Mintz made a cold calculation: sugar supplied the most calories per unit price. For families working twelve-hour factory days with no time to cook, sugar and white bread were the cheapest fuel. A cup of sweet tea welded industrialisation's two engines together: factory gears and plantation rollers. Workers sustained labour with sugar; capital used it to lower wages. Cheap sugar depended on driving labour at the chain's other end to the limits of subsistence.

Mintz offers a sharper distinction between two kinds of meaning. \textbf{Internal meaning} comes from consumers: sugar as reward, comfort or childhood in a birthday cake. \textbf{External meaning} comes from states, capital and empires: tax revenue, shipping routes or reasons to control the Caribbean. For two centuries, British people imagined freely choosing sweetness, while a web far beyond the tea table arranged their choices. Sugar tariffs, naval protection and colonial land systems jointly made sugar cheaper than honey and easier to buy than jam. Taste is never purely private. You think your tongue casts the vote; the ballot has already been printed.

Take the theory's boundary with it. Mintz draws structures well, but individuals can disappear within them: Palmares fugitives, Haitian rebels and boycotting housewives become passive ``structural forces''. Every broad theory risks this: the higher the viewpoint, the smaller the people. Pair it with Equiano's eye. One eye sees the chain from above; the other sees people from the hold.

One final counterexample challenges historical inevitability. Sugar need not grow in slavery's soil. In the early nineteenth century, Napoleon's Continental System cut continental Europe off from colonial cane sugar. France supported beet-sugar extraction; German chemists had identified sugar in beet roots during the eighteenth century. Beet sugar followed another route: European fields, free hired workers, fertiliser and machinery. A substantial part of today's sugar comes from beets. Violence resides not in cane's genes but in institutional calculations. Raw materials do not determine institutions. People choose, or fail to choose, when told ``there is no other way''.

\section{Today's Sugar Bowl}
This three-hundred-year-old problem is in your shopping basket.

Apply the tool today and the sugar story appears in new clothes. Roughly sixty per cent of global cocoa comes from Côte d'Ivoire and Ghana. Research estimates over a million children working in their cocoa plantations, many in conditions approaching forced labour; much of your chocolate begins there. More than half the cobalt in phone batteries comes from the Democratic Republic of Congo, whose mines include children digging by hand. In 2013, Bangladesh's Rana Plaza garment building collapsed, killing over eleven hundred workers making T-shirts for inexpensive Western brands. Workers had been ordered out over wall cracks the day before, then summoned back. Each time, consumers at the chain's end resemble London's little girl: not killers, simply unaware. Each time, ignorance proves part of the chain's design, not an accident.

1791 echoes too. Today's equivalents of sugar-boycotting housewives include fair-trade labels, questions about sweatshop jeans and young people sharing supply-chain investigations. Their effects remain real and limited. Boycotts can raise wrongdoing's costs and force brands to change contracts, but sufficiently long chains and deep poverty let demand find another cheaper corner. Critics of fair-trade certification note that premiums often fail to reach farmers. Critics of boycotts warn that removing children's jobs may push them towards worse livelihoods. Such complexities do not invite a helpless shrug. They remind us that moral impulses need structural knowledge to navigate, or good intentions often miss their target.

\section[Your Turn: Is Sweetness Still Sweet?]{Your Turn: Once You Know, Is Sweetness Still Sweet?}
Return to that four-gram cube.

You now know how a crop moved an empire, how twelve million lives were boiled into cane juice, how abolition mingled conscience, calculation and resistance, and why sweetness was never purely private. Mintz says that eating is never only eating: we consume the final link in an invisible chain.

Here is a question without a standard answer: \textbf{can distance release us from responsibility?}

Make the case for non-involvement fully. Consumers do not perpetrate the violence. The London girl held no whip; you operate no mine. Nobody can investigate every purchase. Requiring everyone to trace every chain effectively outlaws modern life. Now make the case for responsibility. Demand is the market's sole engine; every purchase is a vote. The housewives of 1791 proved informed action possible. Information today is ten thousand times cheaper, making ignorance increasingly a choice rather than a predicament. Neither side is a straw man. Judge for yourself.

A harder layer follows. If responsibility exists, what shape should it take? Boycott, potentially harming equally impoverished workers? Pay more for certified products, potentially enriching certifiers? Promote legislation and customs inspections, slowly and with no guarantee your concern receives attention? Or does merely knowing already change something?

Next time you tear open a sugar packet, consider your answer. It need not match anyone else's, including mine.
